\begingroup
\renewcommand{\thesection}{16A}
\section{Literal History Phase Transport and Its Scope}\label{sec:historyphasegroupoid}
\status{exact replay from supplied native permutations and explicitly transcribed source indices; no decorated history-state identification inferred}
The registered-history phase object provides a concrete test of what reference selection can and cannot settle. To avoid collision with the correlation $\Gamma(A,B)$ or a history denoted $\Gamma$, write its group as $\Gamma_{\rm hist}$. The author-supplied Project~41 excerpts specify a sixteen-element subgroup of the cached $\Aut(G_{60})$ action with
\begin{equation}
\begin{gathered}
\Gamma_{\rm hist}=\langle r,s,z\rangle\cong D_8\times C_2,\\
r^4=s^2=z^2=e,\qquad srs=r^{-1},\qquad [z,r]=[z,s]=e,\\
r=\Aut(G_{60})[3],\quad s=\Aut(G_{60})[325],\quad z=\Aut(G_{60})[324].
\end{gathered}
\tag{HG1}\label{eq:historygrouppresentation}
\end{equation}
The literal cache, rather than a group fingerprint, is used in the new replay. It verifies the presentation, the internal direct product, and generation by $r$ together with the source stabilizer. Completeness of the supplied 480-element automorphism census remains imported; permutation composition and graph preservation are checked directly \citep{cave2026nativepacket,cave2026p4receipts,cave2026draft6replays}.

Set $a=r^2z$, $H_0=\langle a,sz\rangle$, and $K=\langle a\rangle$. The phase set is the left coset space $X=\Gamma_{\rm hist}/H_0$. Its four members and isotropy are
\begin{equation}
S_k=r^kH_0,\quad H_k=r^kH_0r^{-k},\quad
w_k=r^{2(k\bmod2)}sz,\quad H_k=\langle a,w_k\rangle.
\tag{HG2}\label{eq:historyphasecosets}
\end{equation}
The computed $r$ powers have native indices $[0,3,2,5]$, $r^{-1}$ has index $5$, and $K=\{0,326\}$. The native indices are coordinates in the supplied ordered cache, not universal names.

\begin{center}
\small
\begin{tabular}{@{}ccll@{}}
\toprule
$k$ & $r^k$ index & $H_k$ indices & $w_k$ index \\
\midrule
0 & 0 & $\{0,1,323,326\}$ & 1 \\
1 & 3 & $\{0,7,325,326\}$ & 7 \\
2 & 2 & $\{0,1,323,326\}$ & 1 \\
3 & 5 & $\{0,7,325,326\}$ & 7 \\
\bottomrule
\end{tabular}
\end{center}
The phase has period four, while this isotropy and word-generator family has period two. Every $H_k/K$ is a group of order two. Its nonidentity element acts as a history-word exchange in the source interpretation; the quotient group itself is not a torsor with an unnamed identity.

\begin{proposition}[Phase-transport action-groupoid automorphism]\label{prop:phasegroupoid}
For the declared action groupoid $\Gamma_{\rm hist}\ltimes X$, define
\begin{equation}
\alpha(g)=rgr^{-1},\qquad \tau(x)=rx.
\qquad \tau(gx)=\alpha(g)\tau(x).
\tag{HG3}\label{eq:phasegroupoidautomorphism}
\end{equation}
Then $\tau$ on objects and $g\mapsto\alpha(g)$ on arrow labels define an automorphism. It transports $H_k$ to $H_{k+1}$, $S_k$ to $S_{k+1}$ and $w_k$ to $w_{k+1}$, and fixes $K$ setwise.
\end{proposition}
\begin{proof}
Conjugation is a group automorphism, and $r(gx)=(rgr^{-1})(rx)$. Thus an arrow $g:x\to gx$ is sent to $\alpha(g):\tau(x)\to\tau(gx)$. Identities, inverse arrows and composition are preserved; replacing $r$ by $r^{-1}$ gives the inverse functor. Formula (HG2) gives the isotropy and coset identities. The supplied presentation gives $r w_k r^{-1}=w_{k+1}$ and centrality of $a$. All identities are also checked on the literal indexed permutations and all four objects.
\end{proof}
This is the exact finite content of the phase-transport part of Audit018K, sealed in the author's receipt at commit \sourcefile{42ea48f}. Selecting $S_0$ as part of the reference makes $r$ a map to another selected reference, not an automorphism fixing that marked object. Reference covariance is not the assertion that a marking remains pointwise fixed.

\subsection*{Actor transport is not yet a history-state crosswalk}
The reduced continuation interface has two labels $h\in\{2,3\}$, with preserve-roles action $h\mapsto5-h$ and swap-roles action $h\mapsto h$. Their state exchange commutes with both reduced actions. Audit018I records exactly two abstract equivariant bijections to a sign torsor, differing by an overall minus sign. These finite facts are replayed separately.

The stronger inference that (HG3) induces the literal $2\leftrightarrow3$ action requires an additional construction. One needs the appropriate decorated history carrier $\mathcal Y$, its projection $\pi_{\rm hist}:\mathcal Y\to\{2,3\}$ and a lifted automorphism satisfying
\begin{equation}
\pi_{\rm hist}\circ\widetilde\tau=\iota\circ\pi_{\rm hist},
\qquad \iota(2)=3,\quad\iota(3)=2.
\tag{HG4}\label{eq:decoratedhistorytarget}
\end{equation}
The shown Audit018K producer combines the phase theorem with the separately imported reduced-swap result; it does not construct this projection and check (HG4). Conjugating one word-generator index to another transports an actor between fibers. It does not, solely by that computation, determine the permutation of the states in those fibers.

Accordingly this manuscript retains the sealed permutation result, qualifies its stronger state-exchange interpretation, and claims no exhaustive no-go for all richer native history presentations. It does not select a native history-to-central-sign crosswalk or a $\Gamma(A,B)$-sign actor. Program~3 VN03/VN04 remain open at the common-domain action interface. A selected reference can remove a naming ambiguity, but only a construction such as (IR4)--(IR7), or an equivalent native comparison, can establish the coupling.
\endgroup
\setcounter{section}{16}
